% Lesson: Bolzano's theorem (English).

\begin{frame}
\didactatitle{Bolzano's theorem}

\begin{theorem}[Bolzano]
If $f$ is continuous on $[a, b]$ and $f(a)$ and $f(b)$ have opposite signs,
there is $c \in (a, b)$ such that $f(c) = 0$.
\end{theorem}

\bymedium{%
  A continuous function that starts below the axis and ends above it has
  to cross it.
}{%
  A continuous function that starts below the axis and ends above it has
  to cross it. What the theorem adds to the intuition is that this does not
  depend on the picture, but on the real line having no gaps: the
  completeness axiom.
}

\begin{notesonly}
\begin{proof}
Suppose $f(a) < 0 < f(b)$; the other case reduces to this one by replacing $f$
with $-f$. The set $A = \{x \in [a, b] : f(x) < 0\}$ contains $a$ and is
bounded by $b$, so it has a supremum $c \in [a, b]$. We show that $f(c) = 0$.

If $f(c) \neq 0$, taking $\varepsilon = |f(c)|$ in the definition of
continuity there would be a $\delta > 0$ such that $f(x)$ has the same sign as
$f(c)$ for every $x \in [a, b]$ with $|x - c| < \delta$. If $f(c) < 0$,
then $c < b$ and there would be points of $A$ greater than $c$. If $f(c) > 0$,
then $c > a$ and no point of $(c - \delta, c]$ would be in $A$, so
$c - \delta$ would be an upper bound of $A$ smaller than $c$. Both contradict
$c$ being the supremum.

Therefore $f(c) = 0$, and $c \neq a, b$ because neither $f(a)$ nor $f(b)$ vanish.
\end{proof}
\end{notesonly}

\begin{teaching}[To discuss]
What goes wrong if $f$ is not continuous? (The jump $H(x) - \tfrac12$ on $[-1, 1]$.)
And if only
the rationals existed? ($x^2 - 2$ on $[0, 2]$ changes sign and has no
rational root.)
\end{teaching}
\end{frame}

\begin{frame}
\didactatitle{Intermediate values}

\begin{corollary}[Intermediate values]
If $f$ is continuous on $[a, b]$, it takes every value between
$f(a)$ and $f(b)$.
\end{corollary}

\onlynotes{%
  If $k$ lies strictly between $f(a)$ and $f(b)$, it is enough to apply
  Bolzano's theorem to $g(x) = f(x) - k$, which is continuous and changes sign
  on $[a, b]$.
}

\dpause

\begin{example}
The equation $x^3 + x - 1 = 0$ has a solution in $(0, 1)$: the polynomial is
continuous, and it is $-1$ at $0$ and $1$ at $1$.
\onlynotes{%
  Repeating the argument on the half of the interval where the sign changes,
  and then on the half of that half, gives the \keyterm{bisection
  method}: each step halves the length of the interval containing
  the root.
}
\end{example}
\end{frame}
